% 分野別類題: ラプラス方程式（変数分離法）
% 矩形領域・円板領域における u_xx + u_yy = 0 の境界値問題 3 問。
% コンパイル: TEXINPUTS="../../../00_common:$TEXINPUTS" latexmk -xelatex problems.tex
\documentclass[a4paper,11pt]{article}
\usepackage{preamble}

\begin{document}

\kamokumei{類題演習 --- ラプラス方程式（変数分離法）}

\shijibun{次の各境界値問題について、変数分離法を用いて調和関数 $u$ を求めよ（ラプラス方程式 $u_{xx} + u_{yy} = 0$ は、直交座標では $u = X(x)Y(y)$ の積の形を仮定して常微分方程式に帰着でき、極座標では $u = R(r)\Theta(\theta)$ の形を仮定して同様に解ける）。}

\shomon{問1.}{次の矩形領域における境界値問題を解け。}
\[
u_{xx} + u_{yy} = 0 \qquad (0 < x < \pi,\ 0 < y < 1)
\]
\[
u(0, y) = 0, \quad u(\pi, y) = 0, \quad u(x, 0) = 0, \quad u(x, 1) = x(\pi - x)
\]

\shomon{問2.}{次の矩形領域における境界値問題を解け。}
\[
u_{xx} + u_{yy} = 0 \qquad (0 < x < \pi,\ 0 < y < 2)
\]
\[
u(0, y) = 0, \quad u(\pi, y) = 0, \quad u(x, 0) = 0, \quad u(x, 2) = 5\sin 3x
\]

\shomon{問3.}{半径 $3$ の円板 $D = \{(r,\theta) \mid 0 \le r < 3\}$ における次の境界値問題を、極座標形式のラプラス方程式
\[
u_{rr} + \frac{1}{r}u_r + \frac{1}{r^2}u_{\theta\theta} = 0
\]
のもとで解け。ただし $u$ は円板内で有界とする。}
\[
u(3, \theta) = 4 + 6\cos\theta - 3\sin 2\theta \qquad (0 \le \theta < 2\pi)
\]

\end{document}
